\section{The bulk}\label{sec:bulk}

The bulk is the window $|k-pN|=O(B_N)$, where the law of $A_N$ is stable. The
theorem of this section is the whole-row statement for a fresh start. Outside
the bulk it says only that every bin is $o(1/B_N)$, and
Sections~\ref{sec:bigjump} and~\ref{sec:preparation} refine the outside and add
the preparation.

\begin{introtheorem}[the bulk, fresh start]\label{thm:bulk}
Assume (RV) and (SPAN). For each starting sign $\sgn_0$,
\begin{enumerate}[(a)]
\item $\displaystyle\sup_{k\in\Z}\bigl|B_NF^{\sgn_0}_N(k)-h_N(x_{N,k})\bigr|\to0$.
\item $\sup_{N\ge1}\sup_kB_NF^{\sgn_0}_N(k)<\infty$.
\item for every $K>0$, $\displaystyle\sup_{|k-pN|\le KB_N}\Bigl|\frac{B_NF^{\sgn_0}_N(k)}{h_N(x_{N,k})}-1\Bigr|\to0$.
\end{enumerate}
The convergence covers every relation between the two scales $b_{\plus,N}$ and
$b_{\minus,N}$, including ratios that oscillate.
\end{introtheorem}

We prove (a) in Section~\ref{sec:bulk-proof}, and (b) and (c) at its end.

\subsection{Classical inputs}\label{sec:bulk-inputs}

We use four facts about the sums $S^\sgn_m$, stated with the normalization of
Section~\ref{sec:scales}. We write $\bar T_\sgn(x)=\Prob(X_\sgn>x)$, which is
regularly varying with $\bar T_\sgn(x)\sim T_\sgn(x)$ under (RV).

\begin{lemma}[stable and local limits]\label{lem:stable}
Assume (RV) and (SPAN) for the sign $\sgn$, and write $\alpha=\alpha_\sgn$,
$\mu=\mu_\sgn$, $a=a_\sgn$, $g=g_\sgn$, $T=T_\sgn$ and $S_m=S^\sgn_m$. Parts
(i), (iii) and (iv) need only (RV).
\begin{enumerate}[(i)]
\item $Z_m=(S_m-m\mu)/a(m)$ converges in law to $Z_\sgn$, whose density $g$ is
bounded, uniformly continuous, and tends to zero at infinity. Moreover
$m\bigl(\E e^{-(X_\sgn-\mu)/a(m)}-1\bigr)\to K_\alpha$.
\item (Gnedenko's local limit theorem)
$\sup_{k\in\Z}\bigl|a(m)\Prob(S_m=k)-g\bigl((k-m\mu)/a(m)\bigr)\bigr|\to0$ as $m\to\infty$,
and $\bar u_\sgn=\sup_{m\ge1}\sup_ka(m)\Prob(S_m=k)<\infty$.
\item (local large deviations) There is a constant $C_{\mathrm{LD}}=C_{\mathrm{LD}}(\sgn)$ such that
$\Prob(S_m=j)\le C_{\mathrm{LD}}\,m\,T(j-m\mu)/a(m)$
for all $m\ge1$ and all integers $j$ with $j-m\mu\ge1$.
\item (left deviations) The constant $C_-=C_-(\sgn)=\sup_{m\ge1}m\log\E e^{-(X_\sgn-\mu)/a(m)}$
is finite, and
$\Prob(S_m\le m\mu-y)\le e^{C_-}\,e^{-y/a(m)}$ for all $m\ge1$ and $y\ge0$.
\end{enumerate}
\end{lemma}

\begin{proof}
(i) The domain of attraction is classical~\cite[Section~XVII.5]{feller1971}, and we
check the normalization, since the constant $K_\alpha$ is ours. Let $X=X_\sgn$ and, for
$\operatorname{Re}\xi\ge0$, $\Lambda(\xi)=\E[e^{-\xi X}-1+\xi X]$. Since
$e^{-\xi x}-1+\xi x=\int_0^x\xi(1-e^{-\xi y})\,dy$ for $x\ge0$, Fubini's theorem gives,
with $\xi=te^{i\vartheta}$, $t>0$ and $|\vartheta|\le\pi/2$,
\[
 \Lambda(\xi)=\int_0^\infty \xi\bigl(1-e^{-\xi y}\bigr)\bar T(y)\,dy
 =\int_0^\infty e^{i\vartheta}\bigl(1-e^{-e^{i\vartheta}v}\bigr)\bar T(v/t)\,dv .
\]
The integrand is at most $\min(2,v)\bar T(v/t)$ in modulus. As $t\to0$,
$\bar T(v/t)/\bar T(1/t)\to v^{-\alpha}$, and Lemma~\ref{lem:rv}(ii) dominates the ratio
by $2\max(v^{-\alpha-\delta},v^{-\alpha+\delta})$ for $v\ge tx_\delta$, with
$1<\alpha\pm\delta<2$. The part $v<tx_\delta$ contributes at most
$t^2x_\delta^2/\bar T(1/t)\to0$ to $\Lambda(\xi)/\bar T(1/t)$, because $\alpha<2$. Hence,
uniformly in $\vartheta$,
\[
 \Lambda(te^{i\vartheta})=\bar T(1/t)\bigl(F(e^{i\vartheta})+o(1)\bigr),\qquad
 F(\xi)=\int_0^\infty \xi\bigl(1-e^{-\xi v}\bigr)v^{-\alpha}\,dv .
\]
The function $F$ is continuous on $\operatorname{Re}\xi\ge0$, analytic inside, and
$F(t)=t^\alpha F(1)$ for $t>0$. An integration by parts gives
$F(1)=\Gamma(2-\alpha)/(\alpha-1)=K_\alpha$, so $F(\xi)=K_\alpha\xi^\alpha$ with the principal
branch. Next $\E e^{-\xi(X-\mu)}-1=\Lambda(\xi)+O(|\xi|^2)$ as $\xi\to0$. With
$\xi=-i\theta/a(m)$, $m\bar T(a(m)/|\theta|)\to|\theta|^\alpha$ and $m/a(m)^2\to0$ give
$m(\E e^{i\theta(X-\mu)/a(m)}-1)\to K_\alpha(-i\theta)^\alpha$, and L\'evy's continuity
theorem gives $Z_m\to Z_\sgn$. The characteristic function of $Z_\sgn$ has modulus
$\exp(K_\alpha\cos(\pi\alpha/2)|\theta|^\alpha)$, which is integrable, so Fourier inversion
and the Riemann--Lebesgue lemma give the properties of $g$. With $\xi=1/a(m)$ the same
expansion gives the last claim.

(ii) The support of $X_\sgn-\mu$ lies in $\Z-\mu$, and by (SPAN) its maximal lattice
span is $1$. Given (i), Gnedenko's local limit
theorem~\cite[Theorem~4.2.1]{ibragimovlinnik1971} (see
also~\cite[Theorem~8.4.1]{bgt1989}) gives the uniform limit. The norming sequence
there may be any $a_m\sim a(m)$, since $g$ is bounded, uniformly continuous and tends
to zero at infinity. For $\bar u_\sgn$, the limit gives an $m_0$ with
$a(m)\Prob(S_m=k)\le\|g\|_\infty+1$ for $m\ge m_0$, and $a(m)\Prob(S_m=k)\le a(m_0)$ for
$m<m_0$.

(iii) Caravenna and Doney~\cite[Theorem~1.1]{caravennadoney2019} assume right and
left tails $p/A(x)$ and $q/A(x)$ (or $o(1/A(x))$) with $A$ regularly varying of index
in $(0,1)\cup(1,2)$, $p>0$, and mean zero when the index exceeds one. Their bound
(1.6) then holds with no lattice assumption, for all $m\ge1$, $x\ge0$ and bounded
intervals $\mathcal I$, and reads
$\Prob(S_m-m\mu\in x+\mathcal I)\le C_0'(\mathcal I)\,m/(a_mA(x))$ with $a_m=A^{-1}(m)$.
Under (RV) the variable $X_\sgn-\mu$ satisfies these hypotheses with $p=1$, $q=0$ and
$A(x)\sim1/T(x)$. We take $\mathcal I=(-1,0]$ and $x=j-m\mu$, so that
the event contains $\{S_m=j\}$. The asymptotic inverses $a_m$ and $a(m)$ satisfy
$a_m/a(m)\to1$~\cite[Theorem~1.5.12]{bgt1989}, so $a(m)\le c\,a_m$ for all $m$, and
$1/A(x)\le c\,T(x)$ for $x\ge1$. This gives (iii).

(iv) By (i), $m(\E e^{-(X_\sgn-\mu)/a(m)}-1)$ is bounded, and $\log(1+z)\le z$ makes
$C_-$ finite. Markov's inequality for $e^{-(S_m-m\mu)/a(m)}$ and independence give
$\Prob(S_m-m\mu\le-y)\le e^{-y/a(m)}(\E e^{-(X_\sgn-\mu)/a(m)})^m\le e^{C_--y/a(m)}$.
This is the exponential bound behind the Fuk--Nagaev inequalities~\cite{nagaev1979}.
\end{proof}

\subsection{Kernels of one sign}\label{sec:bulk-kernels}

For each sign $\sgn$ the fresh identities use three kernels,
\[
 \mathsf k^{\sgn,1}_m=u^\sgn_m,\qquad \mathsf k^{\sgn,2}_m=q^\sgn_m,\qquad
 \mathsf k^{\sgn,3}_m=q^\sgn_{m+1}\qquad(m\ge1),
\]
with weights $c_{\mathsf k^{\sgn,1}}=1$ and $c_{\mathsf k^{\sgn,2}}=c_{\mathsf k^{\sgn,3}}=\mu_\sgn$. A
kernel of sign $\sgn$ is one of these three sequences, and we write $\mathsf k_m$ and
$c_{\mathsf k}$ when the choice does not matter. We use four facts.
\begin{enumerate}[(K1)]
\item $\sum_{m\ge1}u^\sgn_m(j)\le1$, $\sum_{m\ge1}q^\sgn_m(j)=1$ and
$\sum_{m\ge1}q^\sgn_{m+1}(j)\le1$ for $j\ge1$, by~\eqref{eq:kernel-sums}.
\item $\mathsf k_m(j)=0$ for $m>j$, because $S^\sgn_m\ge m$.
\item Let $J_\sgn(j)=\min\{m\ge1:S^\sgn_m\ge j\}$. Then $\Prob(J_\sgn(j)=m)=q^\sgn_m(j)$
by~\eqref{eq:q-cover}, and $u^\sgn_m(j)\le q^\sgn_m(j)$, since $\{S^\sgn_m=j\}\subset\{J_\sgn(j)=m\}$.
\item $\sum_jq^\sgn_m(j)=\mu_\sgn$ and $\sum_ju^\sgn_m(j)=1$.
\end{enumerate}

\begin{lemma}[kernel local limits]\label{lem:kernels}
Assume (RV) and (SPAN). For each sign $\sgn$ and each kernel $\mathsf k$ of sign $\sgn$,
\[
 \delta_{\mathsf k}(m)=\sup_{j\in\Z}\bigl|a_\sgn(m)\mathsf k_m(j)-c_{\mathsf k}\,g_\sgn\bigl((j-m\mu_\sgn)/a_\sgn(m)\bigr)\bigr|\to0
 \qquad(m\to\infty),
\]
and $\mathsf k_m(j)\le C_3/a_\sgn(m)$ for all $m\ge1$ and $j$, where
$C_3=C_3(\sgn)=\max\bigl(a_\sgn(1),\,\mu_\sgn r_\sgn\bar u_\sgn\bigr)$ and
$r_\sgn=\sup_{m\ge2}a_\sgn(m)/a_\sgn(m-1)<\infty$.
\end{lemma}

\begin{proof}
For $u^\sgn_m$ this is Lemma~\ref{lem:stable}(ii). The kernel $q^\sgn_m$ is a sum of
shifts of $u^\sgn_{m-1}$ with weights $T_\sgn(v)$ of total mass $\mu_\sgn$, and
$q^\sgn_{m+1}$ is the same sum built on $u^\sgn_m$. This gives the bounds, with
$q^\sgn_1=T_\sgn\le1$ at $m=1$. For the limit of $q^\sgn_m$ write $a=a_\sgn$, $g=g_\sgn$,
$z=(j-m\mu)/a(m)$, and fix $v_0\ge1$ and $Z>0$. The shifts $v>v_0$ contribute at most
$(\bar u_\sgn+\|g\|_\infty)\sum_{v>v_0}T(v)$ to $|a(m-1)q_m(j)-\mu g(z)|$. For $v\le v_0$
the point $(j-v-(m-1)\mu)/a(m-1)=z\,a(m)/a(m-1)+(\mu-v)/a(m-1)$ lies within
$\epsilon_m=Z|a(m)/a(m-1)-1|+(v_0+\mu)/a(m-1)$ of $z$ when $|z|\le Z$, and has modulus at
least $Z/2$ for large $m$ when $|z|>Z$. So for large $m$
\[
 \bigl|a(m-1)q_m(j)-\mu g(z)\bigr|\le\mu\bigl(\delta_{u}(m-1)+\omega_g(\epsilon_m)+\bar g(Z/2)\bigr)
 +(\bar u_\sgn+\|g\|_\infty)\sum_{v>v_0}T(v),
\]
where $\omega_g$ is the modulus of continuity of $g$ and $\bar g(r)=\sup_{|y|\ge r}g(y)$.
Let $m\to\infty$, then $v_0,Z\to\infty$, and use $a(m)/a(m-1)\to1$. The kernel
$q^\sgn_{m+1}$ is the same with the shift $-v/a(m)$ in place of $(\mu-v)/a(m-1)$.
\end{proof}

\begin{lemma}[small counts]\label{lem:small}
Assume (RV). For each sign $\sgn$ there is $C_4=C_4(\sgn)$ such that every kernel
of sign $\sgn$ satisfies
\[
 \mathsf k_m(j)\le C_4\bigl(1+m/a_\sgn(m)\bigr)T_\sgn(j)\qquad\text{whenever }8m\mu_\sgn\le j .
\]
Consequently, for $0<\beta\le1$ there is $C_5=C_5(\sgn,\beta)$ with
\begin{equation}\label{eq:bulk-small}
 \mathsf k_m(j)\le C_5\,e_{\sgn,N},\qquad e_{\sgn,N}=\bigl(1+N/a_\sgn(N)\bigr)T_\sgn(N),
\end{equation}
for all $N\ge1$, all $\beta N\le j\le N$ and all $m$ with $8m\mu_\sgn\le\beta N$.
Moreover $B_Ne_{\sgn,N}\to0$.
\end{lemma}

\begin{proof}
Write $T=T_\sgn$, $a=a_\sgn$, $\mu=\mu_\sgn$, and let $C_T$ be the constant of
Lemma~\ref{lem:rv}(ii) with $T(x/2)\le C_TT(x)$. If $4m\mu\le i$, then
$i-m\mu\ge i/2\ge1$, and Lemma~\ref{lem:stable}(iii) gives
\begin{equation}\label{eq:bulk-u-small}
 u_m(i)\le C_{\mathrm{LD}}\,m\,T(i/2)/a(m)\le C_{\mathrm{LD}}C_T\,m\,T(i)/a(m).
\end{equation}
This is the claim for $u_m$. Let $8m\mu\le j$. For $q_1=T$ there is nothing to prove.
For $q_m$ with $m\ge2$ we split $\sum_vT(v)u_{m-1}(j-v)$ at $v=j/2$. If $v\le j/2$, then
$j-v\ge j/2\ge4(m-1)\mu$, so~\eqref{eq:bulk-u-small} bounds $u_{m-1}(j-v)$ by
$C_{\mathrm{LD}}C_T^2r_\sgn\,m\,T(j)/a(m)$, and the weights $T(v)$ add up to at most
$\mu$. If $v>j/2$, then $T(v)\le C_TT(j)$, and the values $u_{m-1}(j-v)$ add up to at
most $1$ by (K4). So $q_m(j)\le(\mu C_{\mathrm{LD}}C_T^2r_\sgn\,m/a(m)+C_T)T(j)$. The
kernel $q_{m+1}$ is the same with $u_m$ in place of $u_{m-1}$. This proves the first
claim with $C_4=C_T+\mu C_{\mathrm{LD}}C_T^2r_\sgn$.

For~\eqref{eq:bulk-small}, $m/a(m)$ has the positive index $1-1/\alpha_\sgn$, so
Lemma~\ref{lem:rv}(iii) gives $m/a(m)\le C_PN/a(N)$ for $m\le N$, and
$T(j)\le T(\beta N)\le C_{T,\beta}T(N)$ by Lemma~\ref{lem:rv}(ii). Finally
$B_NNT_\sgn(N)/a_\sgn(N)$ has the index
\[
 \frac1{\alpha_*}-\frac1{\alpha_\sgn}+1-\alpha_\sgn\le1-\alpha_*<0 ,
\]
because $1/\alpha_*-1/\alpha_\sgn=(\alpha_\sgn-\alpha_*)/(\alpha_*\alpha_\sgn)\le\alpha_\sgn-\alpha_*$,
and $B_NT_\sgn(N)\to0$ by Lemma~\ref{lem:rv}(v).
\end{proof}

The next lemma is the renewal theorem in local form, with the index of the run
that covers a level resolved on its natural scale. For $j\ge1$ and $K>0$ let
\[
 I^\sgn_j(K)=\bigl\{m\ge1:\ |j-m\mu_\sgn|\le K\,a_\sgn(j/\mu_\sgn)\bigr\},\qquad
 y^\sgn_m(j)=\frac{j-m\mu_\sgn}{a_\sgn(j/\mu_\sgn)} .
\]
We call $I^\sgn_j(K)$ the index window. For a modulus $\omega$ (a function with
$\omega(r)\to0$ as $r\to0$) let $\mathcal F(\omega)$ be the set of functions $\psi$ on
$\R$ with $|\psi|\le1$ and $|\psi(y)-\psi(y')|\le\omega(|y-y'|)$.

\begin{lemma}[the index of the covering run]\label{lem:index}
Assume (RV) and (SPAN), and let $\mathsf k$ be a kernel of sign $\sgn$.
\begin{enumerate}[(a)]
\item For every $K\ge2$,
\[
 \limsup_{j\to\infty}\sum_{m\notin I^\sgn_j(K)}\mathsf k_m(j)\le\eta_\sgn(K)
 =\Prob(Z_\sgn\ge K-1)+e^{C_-(\sgn)+1-K/2}.
\]
\item For every modulus $\omega$,
\[
 \sup_{\psi\in\mathcal F(\omega)}\Bigl|\sum_{m\ge1}\mathsf k_m(j)\,\psi\bigl(y^\sgn_m(j)\bigr)
 -\frac{c_{\mathsf k}}{\mu_\sgn}\int_\R g_\sgn(y)\psi(y)\,dy\Bigr|\to0\qquad(j\to\infty).
\]
\end{enumerate}
In particular $\sum_mu^\sgn_m(j)\to1/\mu_\sgn$, which is~\eqref{eq:renewal-thm}.
\end{lemma}

\begin{proof}
Write $\mu=\mu_\sgn$, $a=a_\sgn$, $J=J_\sgn(j)$ and $m_\pm=(j\pm Ka(j/\mu))/\mu$. Note that
$a(x)/x\to0$.

(a) Write $I=I^\sgn_j(K)$. By (K3), $\sum_{m\notin I}u_m(j)$ and $\sum_{m\notin I}q_m(j)$
are at most $\Prob(J\notin I)$, and $\sum_{m\notin I}q_{m+1}(j)\le\Prob(J-1\notin I)$.
All three are at most $\Prob(J\le m_-+1)+\Prob(J>m_+)$. A count of runs below its
typical value forces the sum of the runs above its mean, a right deviation, and a
count above its typical value forces the sum below its mean, a left deviation.

\emph{Too few runs.} Let $m'=\lfloor m_-+1\rfloor$. If $J\le m'$, then $S_{m'}\ge j$. Now
$j-m'\mu\ge Ka(j/\mu)-\mu$, and $a(m')\le a(j/\mu)$ for large $j$. So
\[
 \Prob(J\le m_-+1)\le\Prob\Bigl(\frac{S_{m'}-m'\mu}{a(m')}\ge K-\frac{\mu}{a(j/\mu)}\Bigr)
 \le\Prob(Z_{m'}\ge K-1)
\]
for large $j$. Since $m'\to\infty$, Lemma~\ref{lem:stable}(i) gives
$\limsup_j\Prob(Z_{m'}\ge K-1)\le\Prob(Z_\sgn\ge K-1)$.

\emph{Too many runs.} Let $m''=\lfloor m_+\rfloor$. If $J>m_+$, then $S_{m''}<j$, and
$m''\mu-j\ge Ka(j/\mu)-\mu$. The local bound Lemma~\ref{lem:stable}(iii) does not apply
to this left deviation, and Lemma~\ref{lem:stable}(iv) gives
\[
 \Prob(J>m_+)\le\exp\Bigl(C_--\frac{Ka(j/\mu)-\mu}{a(m'')}\Bigr)\le e^{C_-+1-K/2}
\]
for large $j$, since $m''\le2j/\mu$, $a(2j/\mu)\le2a(j/\mu)$ and $\mu/a(j/\mu)\le1$
eventually.

(b) Fix $K\ge2$. The terms with $m\notin I^\sgn_j(K)$ contribute at most
$\eta_\sgn(K)+o(1)$ by (a), and $\int_{|y|>K}g_\sgn\le\Prob(|Z_\sgn|>K)$. For
$m\in I^\sgn_j(K)$ we have $m\mu/j\to1$ uniformly, so $a(m)/a(j/\mu)\to1$ uniformly.
With $y_m=y^\sgn_m(j)\in[-K,K]$ the argument $(j-m\mu)/a(m)=y_m\,a(j/\mu)/a(m)$ differs
from $y_m$ by at most $K|a(j/\mu)/a(m)-1|$. Lemma~\ref{lem:kernels} and the uniform
continuity of $g_\sgn$ give, uniformly in $m\in I^\sgn_j(K)$,
\[
 \mathsf k_m(j)=\frac{c_{\mathsf k}\,g_\sgn(y_m)+o(1)}{a(j/\mu)} .
\]
The points $y_m$ form an arithmetic progression with step $\Delta_j=\mu/a(j/\mu)\to0$
in $[-K,K]$, with at most $2K/\Delta_j+1$ points, so the $o(1)$ terms add up to $o(1)$.
What remains is $c_{\mathsf k}/\mu$ times a Riemann sum of $g_\sgn\psi$ over $[-K,K]$ with
mesh $\Delta_j$, whose distance from $\int_{-K}^Kg_\sgn\psi$ is at most
$2K(\omega_{g_\sgn}(\Delta_j)+\|g_\sgn\|_\infty\omega(\Delta_j))+2\Delta_j\|g_\sgn\|_\infty$,
uniformly in $\psi\in\mathcal F(\omega)$. Hence the $\limsup$ of the supremum in (b) is
at most $\eta_\sgn(K)+(c_{\mathsf k}/\mu)\Prob(|Z_\sgn|>K)$, and we let $K\to\infty$.
With $\psi\equiv1$ and $\mathsf k=u^\sgn$ we get~\eqref{eq:renewal-thm}.
\end{proof}

\subsection{Proof of Theorem~\ref{thm:bulk}}\label{sec:bulk-proof}

We write $F^{\sgn_0}_N(k)$ through the fresh identity as a sum over the index $m$ of
products of a kernel of sign $\plus$ and a kernel of sign $\minus$. In each product
we resolve the factor of the dominant sign $\tau$ by its local limit and sum the
factor of the other sign through Lemma~\ref{lem:index}. The bins near the ends, the
indices far from the index window of the other sign, the remote bins and the central
bins are treated in turn, and no assumption on the ratio of the two scales enters.

\begin{proof}[Proof of Theorem~\ref{thm:bulk}(a)]
Throughout we fix
\[
 \varepsilon=\min\Bigl\{\frac p2,\ \frac{1-p}2,\ \frac1{16\mu_\plus},\ \frac1{16\mu_\minus}\Bigr\},
\]
and we write $\mu_{\max}=\max_\sgn\mu_\sgn$, $\mu_{\min}=\min_\sgn\mu_\sgn$,
$\eta(K)=\max_\sgn\eta_\sgn(K)$ and $\delta(n)=\sup_{m\ge n}\max_{\sgn,\mathsf k}\delta_{\mathsf k}(m)$,
the maximum over the six kernels. Then $\eta(K)\to0$ as $K\to\infty$ and
$\delta(n)\to0$ as $n\to\infty$. A constant called $C_i$ or $c_i$ depends only on the
laws of the runs and on $\varepsilon$, never on $N$, $k$, $K$ or $R$.

\paragraph{Step 1. The fresh identity as a sum of products.}
We write $\mathcal E_N(k)=|B_NF^{\sgn_0}_N(k)-h_N(x_{N,k})|$ for $k\in\Z$. Let
$1\le k\le N-1$ and write $k_\plus=k$ and $k_\minus=N-k$. By
Proposition~\ref{prop:fresh-identity},
\begin{equation}\label{eq:bulk-pairs}
 F^{\sgn_0}_N(k)=\sum_{i=1}^2\sum_{m\ge1}\mathsf k^{\plus,i}_m(k_\plus)\,\mathsf k^{\minus,i}_m(k_\minus),
\end{equation}
where $(\mathsf k^{\plus,1},\mathsf k^{\minus,1},\mathsf k^{\plus,2},\mathsf k^{\minus,2})$ is
$(q^\plus_{m+1},u^\minus_m,u^\plus_m,q^\minus_m)$ for $\sgn_0=\plus$ and
$(u^\plus_m,q^\minus_{m+1},q^\plus_m,u^\minus_m)$ for $\sgn_0=\minus$. Each is a kernel of
the sign it carries, and in both cases
\begin{equation}\label{eq:bulk-weights}
 c_{\mathsf k^{\plus,1}}c_{\mathsf k^{\minus,1}}+c_{\mathsf k^{\plus,2}}c_{\mathsf k^{\minus,2}}=\mu_\plus+\mu_\minus=s .
\end{equation}
The endpoint bins are $T_\sgn(N)$ or $0$, and $B_NT_\sgn(N)\to0$ by
Lemma~\ref{lem:rv}(v). The profile at $x_{N,k}$ for $k\le0$ or $k\ge N$ is at most
$\varepsilon_h(\min(p,1-p)N/B_N)$, which tends to zero by Lemma~\ref{lem:profile}(b).
So it remains to treat $1\le k\le N-1$.

We also need a uniform bound for one factor. Let $N_1$ be so large that for $N\ge N_1$
and both signs $C_5(\sgn,\varepsilon)b_{\sgn,N}e_{\sgn,N}\le1$ and
$a_\sgn(N/s)\le(8/\varepsilon)a_\sgn(\varepsilon N/(8\mu_\sgn))$. The second inequality
holds eventually because the ratio tends to
$(8\mu_\sgn/(\varepsilon s))^{1/\alpha_\sgn}<8/\varepsilon$. Then for $N\ge N_1$,
$\varepsilon N\le j\le N$, all $m\ge1$ and every kernel $\mathsf k$ of sign $\sgn$,
\begin{equation}\label{eq:bulk-factor}
 \mathsf k_m(j)\le\frac{C_6}{b_{\sgn,N}},\qquad C_6=1+8C_3/\varepsilon .
\end{equation}
Indeed, for $8m\mu_\sgn\le\varepsilon N$ this is~\eqref{eq:bulk-small} with
$\beta=\varepsilon$, and for larger $m$ Lemma~\ref{lem:kernels} gives
$\mathsf k_m(j)\le C_3/a_\sgn(\varepsilon N/(8\mu_\sgn))\le 8C_3/(\varepsilon b_{\sgn,N})$. For
the dominant sign the bound is $C_6/B_N$. For the other sign it may be much larger,
and we never use it there. In every product of~\eqref{eq:bulk-pairs} we bound the
factor of the dominant sign and sum the factor of the other sign.

\paragraph{Step 2. Small counts and the ends of the row.}
Let $k_\minus<\varepsilon N$. By (K2) every nonzero product in~\eqref{eq:bulk-pairs} has
$m\le k_\minus<\varepsilon N$. The plus factor sits at $k_\plus>(1-\varepsilon)N\ge N/2$,
and $8m\mu_\plus\le N/2$, so Lemma~\ref{lem:small} with $\beta=1/2$ bounds it by
$C_5(\plus,\tfrac12)e_{\plus,N}$. The minus factors add up to at most $1$ by (K1).
Hence $B_NF^{\sgn_0}_N(k)\le2C_5(\plus,\tfrac12)B_Ne_{\plus,N}\to0$ uniformly in these $k$,
and the case $k_\plus<\varepsilon N$ is the same with the signs exchanged. In both cases
$|k-pN|\ge\varepsilon N$, so the profile is at most $\varepsilon_h(\varepsilon N/B_N)\to0$.
So
\begin{equation}\label{eq:bulk-ends}
 \sup_{\min(k,N-k)<\varepsilon N}\mathcal E_N(k)\to0 .
\end{equation}
From now on $\varepsilon N\le k_\plus,k_\minus\le N$, and~\eqref{eq:bulk-factor} applies
to both factors.

\paragraph{Step 3. Indices far from the index window of the other sign.}
Fix $K\ge2$ and write $I=I^{\bar\tau}_{k_{\bar\tau}}(K)$. By~\eqref{eq:bulk-factor} for
the dominant factor and Lemma~\ref{lem:index}(a) for the other factor at the level
$k_{\bar\tau}\ge\varepsilon N$,
\begin{equation}\label{eq:bulk-remote-counts}
 \limsup_{N\to\infty}\ \sup_{\varepsilon N\le k\le(1-\varepsilon)N}\
 B_N\sum_{i=1}^2\sum_{m\notin I}\mathsf k^{\plus,i}_m(k_\plus)\mathsf k^{\minus,i}_m(k_\minus)\le2C_6\,\eta(K).
\end{equation}
For $m\in I$ we record three facts. Let $N\ge N_2(K)$, where $N_2(K)\ge N_1$ is so
large that $Ka_\sgn(N/\mu_\sgn)\le\varepsilon N/2$ for both signs. Then every $m\in I$
satisfies
\begin{equation}\label{eq:bulk-window}
 \frac{\varepsilon N}{2\mu_{\bar\tau}}\le m\le\frac{2N}{\mu_{\bar\tau}} ,
\end{equation}
so Lemma~\ref{lem:kernels} applies to the dominant factor with the error
$\delta(\lceil\varepsilon N/(2\mu_{\max})\rceil)\to0$. Next, the definitions of $p$ and
$s$ give the exact identity
\begin{equation}\label{eq:bulk-shift}
 k_\tau-m\mu_\tau=\varsigma_\tau\frac{s(k-pN)}{\mu_{\bar\tau}}
 +\frac{\mu_\tau}{\mu_{\bar\tau}}\bigl(k_{\bar\tau}-m\mu_{\bar\tau}\bigr).
\end{equation}
For $\tau=\plus$ the right side is
$(sk-\mu_\plus N+\mu_\plus k_\minus)/\mu_\minus-m\mu_\plus=k-m\mu_\plus$, and $\tau=\minus$ is the
same computation. Finally, by Lemma~\ref{lem:rv}(ii) there is $N_3\ge N_2(K)$ such that
for $N\ge N_3$
\begin{equation}\label{eq:bulk-ratios}
 a_\sgn(N/\mu_\sgn)\le c_0\,b_{\sgn,N},\qquad
 a_\sgn(2N/\mu_{\min})\le c_1b_{\sgn,N},\qquad
 B_N\le c_2\,a_\tau(\varepsilon N/(2\mu_{\max})),
\end{equation}
with $c_0=s/\mu_{\min}$, $c_1=2s/\mu_{\min}$ and $c_2=2/\varepsilon$, because the ratios
tend to $(s/\mu_\sgn)^{1/\alpha_\sgn}$, $(2s/\mu_{\min})^{1/\alpha_\sgn}$ and
$(2\mu_{\max}/(\varepsilon s))^{1/\alpha_\tau}$.

\paragraph{Step 4. Remote bins.}
Let $R\ge4\mu_{\max}/\mu_{\min}$ and put $K=K(R)=sR/(2c_0\mu_{\max})=\mu_{\min}R/(2\mu_{\max})\ge2$.
Let $N\ge N_3$ with $N_3$ taken for this $K$, and let $\varepsilon N\le k\le(1-\varepsilon)N$
and $|k-pN|>RB_N$. For $m\in I$ we have
$|k_{\bar\tau}-m\mu_{\bar\tau}|\le Ka_{\bar\tau}(N/\mu_{\bar\tau})\le c_0KB_N$
by~\eqref{eq:bulk-ratios}, so~\eqref{eq:bulk-shift} gives
\[
 |k_\tau-m\mu_\tau|\ge\frac{sRB_N}{\mu_{\bar\tau}}-\frac{\mu_\tau}{\mu_{\bar\tau}}c_0KB_N
 \ge\frac{sRB_N}{2\mu_{\bar\tau}} .
\]
By~\eqref{eq:bulk-window} and~\eqref{eq:bulk-ratios}, $a_\tau(m)\le c_1B_N$, so the scaled
argument $z=(k_\tau-m\mu_\tau)/a_\tau(m)$ of the dominant factor has modulus at least
$c_8R$ with $c_8=s/(2c_1\mu_{\max})$. Lemma~\ref{lem:kernels} and~\eqref{eq:bulk-ratios}
give
\[
 B_N\mathsf k^{\tau,i}_m(k_\tau)\le\frac{B_N}{a_\tau(m)}\bigl(\mu_{\max}\bar g_\tau(c_8R)+\delta(m)\bigr)
 \le c_2\bigl(\mu_{\max}\bar g_\tau(c_8R)+\delta(\lceil\varepsilon N/(2\mu_{\max})\rceil)\bigr).
\]
The factors of sign $\bar\tau$ add up to at most $1$ for each $i$ by (K1). Together
with~\eqref{eq:bulk-remote-counts} and Lemma~\ref{lem:profile}(b),
\begin{equation}\label{eq:bulk-remote-bins}
 \limsup_{N\to\infty}\sup_{\substack{\varepsilon N\le k\le(1-\varepsilon)N\\|k-pN|>RB_N}}\mathcal E_N(k)
 \le e(R)=2c_2\mu_{\max}\max_\sgn\bar g_\sgn(c_8R)+2C_6\eta(K(R))+\varepsilon_h(R),
\end{equation}
and $e(R)\to0$ as $R\to\infty$.

\paragraph{Step 5. The central sum.}
Fix $R\ge1$ and $K\ge2$, and let $|k-pN|\le RB_N$, so $x=x_{N,k}\in[-R,R]$. For large
$N$ these bins satisfy $\varepsilon N\le k\le(1-\varepsilon)N$, since $B_N=o(N)$. Let
$m\in I$ and $y_m=y^{\bar\tau}_m(k_{\bar\tau})\in[-K,K]$. Since
$k_\plus/\mu_\plus=N/s+xB_N/\mu_\plus$ and $k_\minus/\mu_\minus=N/s-xB_N/\mu_\minus$, we have
$k_{\bar\tau}/\mu_{\bar\tau}=(N/s)(1+o(1))$ and $m=(N/s)(1+o(1))$, uniformly in $|x|\le R$
and $m\in I$. Lemma~\ref{lem:rv}(i), applied to $a_\plus$ and to $a_\minus$ separately,
gives
\[
 a_{\bar\tau}(k_{\bar\tau}/\mu_{\bar\tau})=b_{\bar\tau,N}(1+o(1)),\qquad a_\tau(m)=B_N(1+o(1)),
\]
uniformly in $|x|\le R$ and $m\in I$, whichever sign is dominant. Since
$b_{\bar\tau,N}=\theta_{\bar\tau,N}B_N$, the identity~\eqref{eq:bulk-shift} gives for the
argument of the dominant factor
\[
 z_m=\frac{k_\tau-m\mu_\tau}{a_\tau(m)}
 =\varsigma_\tau\frac{sx}{\mu_{\bar\tau}}\,r_{1}
 +\frac{\mu_\tau}{\mu_{\bar\tau}}\theta_{\bar\tau,N}\,y_m\,r_{2},\qquad
 r_1=\frac{B_N}{a_\tau(m)},\quad r_2=\frac{a_{\bar\tau}(k_{\bar\tau}/\mu_{\bar\tau})}{b_{\bar\tau,N}}\,\frac{B_N}{a_\tau(m)} ,
\]
where $r_1,r_2\to1$ uniformly. Since $|x|\le R$, $|y_m|\le K$ and $\theta_{\bar\tau,N}\le1$,
\[
 \sup_{|x|\le R}\ \sup_{m\in I}\ \bigl|z_m-\zeta_{N,x}(y_m)\bigr|\to0,\qquad
 \zeta_{N,x}(y)=\varsigma_\tau\frac{sx}{\mu_{\bar\tau}}+\frac{\mu_\tau}{\mu_{\bar\tau}}\theta_{\bar\tau,N}\,y .
\]
Lemma~\ref{lem:kernels}, $B_N/a_\tau(m)\to1$ and the uniform continuity of $g_\tau$ give,
uniformly in $|x|\le R$ and $m\in I$,
\begin{equation}\label{eq:bulk-dominant}
 B_N\mathsf k^{\tau,i}_m(k_\tau)=c_{\mathsf k^{\tau,i}}\,g_\tau\bigl(\zeta_{N,x}(y_m)\bigr)+o(1).
\end{equation}
Let $\psi_{N,x}=g_\tau\circ\zeta_{N,x}/\|g_\tau\|_\infty$. The slope of $\zeta_{N,x}$ is at most
$\mu_{\max}/\mu_{\min}$, so $\psi_{N,x}\in\mathcal F(\omega)$ for all $N$ and $x$, with
$\omega(r)=\max_\sgn\omega_{g_\sgn}(\mu_{\max}r/\mu_{\min})/\|g_\sgn\|_\infty$. We
multiply~\eqref{eq:bulk-dominant} by $\mathsf k^{\bar\tau,i}_m(k_{\bar\tau})$ and sum over
$m\in I$. The $o(1)$ terms add up to $o(1)$ by (K1). Then we add back the indices
$m\notin I$ in the sum of $\mathsf k^{\bar\tau,i}_m\psi_{N,x}(y_m)$, at a cost of at most
$\eta(K)+o(1)$ by Lemma~\ref{lem:index}(a), and apply Lemma~\ref{lem:index}(b) at the
level $k_{\bar\tau}\ge\varepsilon N$. This gives, uniformly in $|x|\le R$,
\[
 B_N\sum_{m\in I}\mathsf k^{\plus,i}_m(k_\plus)\mathsf k^{\minus,i}_m(k_\minus)
 =\frac{c_{\mathsf k^{\tau,i}}c_{\mathsf k^{\bar\tau,i}}}{\mu_{\bar\tau}}\int_\R g_{\bar\tau}(y)\,g_\tau\bigl(\zeta_{N,x}(y)\bigr)dy
 +O\bigl(\eta(K)\bigr)+o(1),
\]
with the constant $\mu_{\max}\max_\sgn\|g_\sgn\|_\infty$ in $O(\eta(K))$. We sum over $i$
and use~\eqref{eq:bulk-weights}. By Lemma~\ref{lem:profile}(e) the integral with the
weight $s/\mu_{\bar\tau}$ is exactly $h_N(x)$. Adding~\eqref{eq:bulk-remote-counts} for
the indices outside $I$,
\begin{equation}\label{eq:bulk-central}
 \limsup_{N\to\infty}\sup_{|k-pN|\le RB_N}\mathcal E_N(k)\le C_7\,\eta(K),
\end{equation}
with $C_7=2C_6+2\mu_{\max}\max_\sgn\|g_\sgn\|_\infty$. Since $K$ is arbitrary, the left
side is zero.

\paragraph{Combination.} Let $R\ge4\mu_{\max}/\mu_{\min}$. Every $k\in\Z$ is an endpoint
or lies outside the row (Step 1), or within $\varepsilon N$ of an end (Step 2), or is a
remote bin with $|k-pN|>RB_N$ (Step 4), or a central bin with $|k-pN|\le RB_N$
(Step 5). By~\eqref{eq:bulk-ends},~\eqref{eq:bulk-remote-bins}
and~\eqref{eq:bulk-central}, $\limsup_N\sup_{k\in\Z}\mathcal E_N(k)\le e(R)$, and we let
$R\to\infty$. No step assumed that $\theta_N$ converges, since every error term is a
maximum over the two signs of quantities that do not involve $\theta_N$, and
$\theta_N$ enters only through the slope of $\zeta_{N,x}$, which is at most
$\mu_{\max}/\mu_{\min}$.
\end{proof}

\begin{proof}[Proof of Theorem~\ref{thm:bulk}(b) and (c)]
For (b), part (a) and Lemma~\ref{lem:profile}(a) give an $N_4$ with
$B_NF^{\sgn_0}_N(k)\le C_1+1$ for all $N\ge N_4$ and all $k$, and
$B_NF^{\sgn_0}_N(k)\le B_N\le B_{N_4}$ for $N<N_4$, since $B_N$ is nondecreasing. For
(c), Lemma~\ref{lem:profile}(d) gives $c_K>0$ with $h_N\ge c_K$ on $[-K,K]$ for all $N$,
and dividing the uniform error of (a) by $c_K$ gives (c).
\end{proof}

The proof used (SPAN) only through Lemma~\ref{lem:stable}(ii), and it did not use (LR).
